% 原书PDF第11—21页，4.1.3.6全节；止于4.1.3.7标题前。
\subsubsection{异构体表达定量}
\bookdeepheading{paragraph}{4.1.3.6.1}{RNA-seq 数据无缺口比对读段的生成模型}

表达定量的一项主要挑战，是 RNA-seq 读段常可比对到多个基因或异构体（Li and Dewey 2011）。RNA-seq 基因表达定量的目的，是利用 RNA-seq 数据测量样本中转录本的拷贝数。为准确量化表达，必须考虑读段比对的不确定性、测序的非均匀性，以及多个异构体的存在（Li et al. 2010）。当存在可变剪接时，不同外显子可能被不同数量的异构体共享；许多读段难以唯一比对到其来源异构体（Liu et al. 2015）。因此，需要对比对到共享外显子的读段进行解卷积。潜变量生成模型是常用的读段解卷积方法。

图4.5展示了基因结构与异构体。示例基因包含三个外显子和三个转录本（异构体），图中还给出一个包含全部已注释外显子及连接位点的附加转录本。假设抽取了 $N$ 条独立同分布、长度为 $L$ 的读段。读段序列是观测数据，第 $n$ 条读段的序列用随机变量 $R_n$ 表示。用有向图模型或贝叶斯网络（Li et al. 2010）描述读段生成过程，如图4.6所示。

\begin{figure}[htbp]
\centering
\begin{tikzpicture}[x=1cm,y=1cm,>=Stealth,font=\small]

\draw[red,very thick] (1.5,5)--(4.5,5);
\draw[inkblue,very thick] (6,5)--(9,5);
\draw[orange,very thick] (11.5,5)--(14.5,5);
\node at (3,4.5) {外显子1};\node at (7.5,4.5) {外显子2};\node at (13,4.5) {外显子3};
\foreach \x/\y in {1.5/5.15,2/5.35,2.7/5.55,6/5.3,6.7/5.6,7.5/5.9}
 \draw[green!65!black,thick] (\x,\y)--++(1.5,0);
\node at (2.1,5.9) {读段1};\node at (3.7,5.9) {读段2};\node at (6.5,6.2) {读段3};
\draw[magenta,|<->|] (8.25,6)--(9.3,6);\node at (8.78,6.2) {$L$};
\node[anchor=east] at (1.3,3.4) {转录本1};
\draw[red,very thick] (1.5,3.4)--(4.5,3.4);\draw[inkblue,very thick] (4.5,3.4)--(7.5,3.4);
\node[anchor=east] at (1.3,2.6) {转录本2};
\draw[red,very thick] (1.5,2.6)--(4.5,2.6);\draw[orange,very thick] (4.5,2.6)--(7.5,2.6);
\node[anchor=east] at (1.3,1.8) {转录本3};
\draw[inkblue,very thick] (1.5,1.8)--(4.5,1.8);\draw[orange,very thick] (4.5,1.8)--(7.5,1.8);
\draw[red,very thick] (1.5,.8)--(4.5,.8);
\draw[inkblue,very thick] (4.5,.8)--(7.5,.8);
\draw[orange,very thick] (7.5,.8)--(10.5,.8);
\node[anchor=west] at (10.8,.8) {附加转录本};

\end{tikzpicture}
\origfig{5}{基因结构及异构体示例。}
\label{fig:ch04a-05}
\end{figure}

\begin{figure}[htbp]
\centering
\begin{tikzpicture}[x=1cm,y=1cm,>=Stealth,font=\small]

\node[ellipse,draw=inkblue,minimum width=3cm,minimum height=1.6cm] (ch04agG) at (2,1.9) {};
\node at (2,1.4) {$G_n$};\node[text=red] at (.85,2.2) {$M$};
\draw[magenta] (1.2,2.5)--(2.1,2.5);\draw[green!65!black] (2.1,2.5)--(2.8,2.5);
\draw[magenta] (1.2,2.25)--(2.1,2.25);\draw[inkblue] (2.1,2.25)--(2.5,2.25);
\draw[magenta] (1.2,2)--(2.1,2);\draw[inkblue] (2.1,2)--(3.15,2);
\node at (-.65,1.9) {$\theta$};\draw[green!65!black,->] (-.35,1.9)--(ch04agG.west);
\node[ellipse,draw=inkblue,minimum width=3cm,minimum height=1.65cm] (ch04agS) at (5.3,4.1) {};
\node at (5.3,3.72) {$S_n$};
\draw[inkblue] (4.3,4.4)--(6.3,4.4);
\foreach \x in {4.3,4.75,5.2,5.65,6.1}{\fill[red] (\x,4.4) circle(.06);}
\draw[orange,->] (4.65,3.95)--(4.75,4.35);
\node[ellipse,draw=inkblue,minimum width=2.7cm,minimum height=1.4cm] (ch04agO) at (5.3,.6) {$O_n$};
\node[ellipse,draw=inkblue,minimum width=3.4cm,minimum height=1.9cm] (ch04agR) at (10,2.1) {};
\node at (10,1.55) {$R_n$};\draw[inkblue] (8.95,2.25)--(11.05,2.4);
\foreach \x in {9,9.45,9.9,10.35,10.8}{\fill[red] (\x,2.45) circle(.06);}
\node at (10.65,4.05) {读段序列};\draw[green!65!black,->,thin] (9.9,2.45)--(10.6,3.8);
\node at (1.2,4.7) {$N$};
\draw[green!65!black,->] (ch04agG)--(ch04agS);
\draw[green!65!black,->] (ch04agG)--(ch04agO);
\draw[green!65!black,->] (ch04agG)--(ch04agR);
\draw[green!65!black,->] (ch04agS)--(ch04agR);
\draw[green!65!black,->] (ch04agO)--(ch04agR);

\end{tikzpicture}
\origfig{6}{RNA-seq 数据的图形生成模型。}
\label{fig:ch04a-06}
\end{figure}


考虑 $M$ 个异构体和参数向量 $\theta=[\theta_0,\theta_1,\ldots,\theta_M]$。其中，$\theta_0$ 对应噪声异构体，$\theta_i$（$1\le i\le M$）对应第 $i$ 个异构体。为生成观测读段序列 $R_n$，引入三个隐藏变量：异构体变量 $G_n$ 表示读段 $n$ 来自哪个异构体；起始位置变量 $S_n$ 表示读段 $n$ 的起始位置；方向变量 $O_n$ 表示测序的是哪一条链。在图4.6的贝叶斯网络假设下，变量 $G_n$、$S_n$、$O_n$ 与 $R_n$ 的联合概率为
\begin{equation}
P(g,s,o,r\mid\theta)=\prod_{n=1}^{N}P(g_n\mid\theta)P(s_n\mid g_n)P(o_n\mid g_n)P(r_n\mid g_n,s_n,o_n).
\tag{4.18}
\end{equation}

下面分别计算式 (4.18) 中的各项。异构体随机变量取值范围为 $[0,M]$。设第 $i$ 个异构体在异构体总体中所占比例为 $\theta_i$，则选择第 $i$ 个异构体的概率为 $\theta_i$，即 $P(G_n=i\mid\theta)=\theta_i$，且 $\sum_{i=0}^{M}\theta_i=1$（图4.7a）。设 $l_i$（$i=0,\ldots,M$）为第 $i$ 个异构体的长度。读段起始位置的取值区间为 $[1,\max_i l_i]$。假设读段在转录本上均匀起始，并且可以从异构体的最后一个位置开始、延伸至 poly(A) 尾。在这种情况下，读段沿第 $i$ 个转录本从位置 $j$ 开始生成的概率为
\[
P(S_n=j\mid G_n=i)=\frac1{l_i}
\]
（图4.7b）。若假设 mRNA 末端没有 poly(A) 尾，则
\[
P(S_n=j\mid G_n=i)=\frac1{l_i-L+1}.
\]

\begin{figure}[htbp]
\centering
\begin{tikzpicture}[x=1cm,y=1cm,>=Stealth,font=\small]

\node[anchor=east] at (2.6,8.6) {转录本1};
\draw[red,thick] (2.8,8.6)--(5.8,8.6);\draw[inkblue,thick] (5.8,8.6)--(8,8.6);\node at (9.5,8.6) {$\theta_1$};
\node[anchor=east] at (2.6,7.8) {转录本 $i$};
\draw[red,thick] (2.8,7.8)--(5.8,7.8);\draw[orange,thick] (5.8,7.8)--(9,7.8);\node at (9.5,7.8) {$\theta_i$};
\node[anchor=east] at (2.6,7) {转录本 $M$};
\draw[inkblue,thick] (2.8,7)--(5.8,7);\draw[orange,thick] (5.8,7)--(9,7);\node at (9.5,7) {$\theta_M$};
\foreach \y in {8.4,8.22,8.04,7.6,7.42,7.24}{\draw[inkblue!55,dashed,thin] (2.8,\y)--(9,\y);}
\node[anchor=east] at (2.6,6.25) {附加转录本};
\draw[red,thick] (2.8,6.25)--(5.8,6.25);\draw[inkblue,thick] (5.8,6.25)--(7.6,6.25);\draw[orange,thick] (7.6,6.25)--(10.5,6.25);
\node at (11,6.25) {$\theta_0$};\node at (.6,6) {(a)};
\draw[red,thick] (2.8,4.9)--(5.8,4.9);\draw[orange,thick] (5.8,4.9)--(9,4.9);
\foreach \x in {2.8,3.8,4.8,5.8,6.7,7.6,8.3,9}{\draw[green!65!black] (\x,4.9)--(\x,5.35);}
\draw[inkblue,thick,->] (4.8,4.9)--(4.8,5.35)--(5.6,5.35);
\node at (9.6,5.05) {$l_i$};\node at (5.5,4.35) {$p=\dfrac1{l_i}$};
\draw[red,thick] (2.8,3)--(5.8,3);\draw[orange,thick] (5.8,3)--(7.6,3);\draw[magenta,thick] (7.6,3)--(9,3);
\foreach \x in {2.8,3.8,4.8,5.8,6.7,7.6,8.3,9}{\draw[green!65!black] (\x,3)--(\x,3.45);}
\node at (7.95,3.22) {A};\node at (8.65,3.22) {A};
\draw[inkblue,thick] (2.8,3)--(2.8,2.35);\draw[inkblue,thick] (7.6,3)--(7.6,2.35);
\draw[red,dashdotted] (2.8,2.6)--(7.6,2.6);
\node[fill=white,inner sep=1pt] at (5.2,2.6) {$L_i-L+1$};
\node at (8.3,2.7) {$L$};\node at (11,2.95) {$p=\dfrac1{l_i-L+1}$};\node at (.6,2.2) {(b)};
\node at (2,.8) {$G_n$};\draw[green!65!black,->] (2.3,.8)--(3,.8);
\fill[red!90] (3,.6) rectangle (7.4,1);\fill[inkblue!70] (7.4,.6) rectangle (9.7,1);
\node[text=white,font=\ttfamily\small,inner sep=0] at (5.2,.8) {CTGAGGTTCCCAAGGA};
\node[text=white,font=\ttfamily\small,inner sep=0] at (8.55,.8) {AGGTTCA};
\draw[green!65!black] (5.4,1.05)--(5.4,1.6);
\node[anchor=west] at (5.5,1.5) {起始位置 $j$};
\node[anchor=west] at (10,.9) {异构体 $i$};
\node[anchor=west] at (10,.35) {$w_t(r_{nt},g_{j+t-1}^i)$};
\node at (4.3,-.05) {$R_n$};\draw[green!65!black,->] (4.6,-.05)--(5.3,-.05);
\fill[magenta!75] (5.3,-.3) rectangle (8,-.0);
\node[text=white,font=\ttfamily\small] at (6.65,-.15) {CTGAGGTT};\node at (.6,-.35) {(c)};

\end{tikzpicture}
\origfig{7}{生成模型示意。(a) 选择第 $i$ 种异构体的概率；(b) 沿第 $i$ 个转录本从位置 $j$ 开始生成读段的概率；(c) 给定所选异构体与起始位置，观测到该读段序列的概率。}
\label{fig:ch04a-07}
\end{figure}


RNA-seq 有两种方案：标准的非链特异 RNA-seq，以及链特异 RNA-seq。标准 RNA-seq 方案会丢失原始 mRNA 测序模板来自哪条链的信息；链特异 RNA-seq 的序列读段从第一条链生成，因而在整个测序过程中保留链信息（Zhao et al. 2015）。方向变量 $O_n$ 是二元变量。设 $O_n=0$ 表示读段 $n$ 的序列方向与其来源异构体相同；$O_n=1$ 表示该序列为反向互补序列。因此，对于链特异方案，假设 $O_n=0$，即 $P(O_n=0\mid G_n\ne0)=1$。对于非链特异 RNA-seq，没有可用的链信息，方向变量 $O_n$ 以相等概率取0或1，即
\[
P(O_n=0\mid G_n\ne0)=P(O_n=1\mid G_n\ne0)=0.5.
\]

最后，计算给定所选异构体与起始位置时，观测到读段序列的概率
\[
P(R_n=r_n\mid G_n=i,S_n=j,O_n=k).
\]
该概率根据观测序列与异构体的比较，以及二者的比对评分计算（图4.7c）。定义指示随机变量
\begin{equation}
Z_{nijk}=\begin{cases}
1,&G_n=i,\ S_n=j,\ O_n=k,\\
0,&\text{其他情形}.
\end{cases}
\tag{4.19}
\end{equation}
于是，上述条件概率可写成 $P(R_n=r_n\mid Z_{nijk}=1)$，并计算为
\begin{equation}
P(R_n=r_n\mid Z_{nijk}=1)=\begin{cases}
\displaystyle\prod_{t=1}^{L}w_t(r_{nt},g_{j+t-1}^{i}),&k=0,\\[1ex]
\displaystyle\prod_{t=1}^{L}w_t(r_{nt},\bar g_{j+t-1}^{i}),&k=1,
\end{cases}
\tag{4.20}
\end{equation}
其中 $w_t(a,b)$ 是位置特异的替换矩阵；$r_{nt}$ 为第 $n$ 条读段的第 $t$ 个碱基；$g_{j+t-1}^{i}$ 为第 $i$ 个异构体的第 $j+t-1$ 个碱基；$\bar g_{j+t-1}^{i}g_{j+t-1}^{i}$ 表示该异构体反向互补序列的第 $j+t-1$ 个碱基。替换矩阵为每一对碱基的匹配或错配（替换）赋予评分。$w_t(a,b)$ 的值是在位置 $t$ 将 $a$ 与 $b$ 比对的得分，由位置特异替换矩阵确定。该矩阵或评分不仅考虑匹配／错配，还纳入碱基识别错误、多态性和参考序列错误。这些因素都会导致异构体序列与观测读段之间出现可观察到的替换。

如果读段来自噪声异构体，则假设它由与位置无关的背景分布 $\beta$ 生成。令 $G_n=0$、起始位置 $j=1$、方向变量 $O_n=0$，式 (4.20) 化为
\begin{equation}
P(R_n=r_n\mid Z_{n010}=1)=\prod_{t=1}^{L}\beta(r_{nt}).
\tag{4.21}
\end{equation}

\bookdeepheading{paragraph}{4.1.3.6.2}{RNA-seq 数据有缺口比对读段的生成模型}

上一生成模型的局限，是无法处理读段的有缺口比对。为克服这一局限，对模型作修改，允许读段采用有缺口比对（Nariai et al. 2013）。将图4.6的无缺口比对模型扩展为图4.8的有缺口比对模型，增加两个节点：$A_n$ 处理比对状态，$Q_n$ 纳入测序质量评分。以 $G_n$、$S_n$、$O_n$、$A_n$、$Q_n$ 和 $R_n$ 分别表示转录本异构体选择、转录本起始位置、读段方向、比对状态、序列质量评分，以及读段变量。参数 $\theta$ 的定义见第4.1.3.6.1节。

比对状态变量 $A_n$ 表示读段 $n$ 与参考异构体 $i$ 在每个比对位置上的状态，即匹配／错配及插入缺失。序列质量评分变量 $Q_n$ 取 Phred 质量评分（Nariai et al. 2013）。图4.8有向图模型中各随机变量的联合概率为
\begin{equation}
\begin{aligned}
P(G_n,S_n,O_n,A_n,Q_n,R_n\mid\theta)
={}&P(Q_n)P(G_n\mid\theta)P(S_n\mid G_n)P(O_n\mid G_n)\\
&\times P(A_n\mid G_n,S_n,O_n)P(R_n\mid G_n,O_n,S_n,A_n,Q_n).
\end{aligned}
\tag{4.22}
\end{equation}
\begin{figure}[htbp]
\centering
\begin{tikzpicture}[x=1cm,y=1cm,>=Stealth,font=\small]

\node[ellipse,draw=inkblue,minimum width=3cm,minimum height=1.6cm] (ch04agG) at (2,1.9) {};
\node at (2,1.4) {$G_n$};\node[text=red] at (.85,2.2) {$M$};
\draw[magenta] (1.2,2.5)--(2.1,2.5);\draw[green!65!black] (2.1,2.5)--(2.8,2.5);
\draw[magenta] (1.2,2.25)--(2.1,2.25);\draw[inkblue] (2.1,2.25)--(2.5,2.25);
\draw[magenta] (1.2,2)--(2.1,2);\draw[inkblue] (2.1,2)--(3.15,2);
\node at (-.65,1.9) {$\theta$};\draw[green!65!black,->] (-.35,1.9)--(ch04agG.west);
\node[ellipse,draw=inkblue,minimum width=3cm,minimum height=1.65cm] (ch04agS) at (3.5,4.9) {};
\node at (3.5,4.52) {$S_n$};
\draw[inkblue] (2.5,5.2)--(4.5,5.2);
\foreach \x in {2.5,2.95,3.4,3.85,4.3}{\fill[red] (\x,5.2) circle(.06);}
\draw[orange,->] (2.85,4.75)--(2.95,5.15);
\node[ellipse,draw=inkblue,minimum width=2.7cm,minimum height=1.4cm] (ch04agO) at (5.3,.6) {$O_n$};
\node[ellipse,draw=inkblue,minimum width=3.4cm,minimum height=1.9cm] (ch04agR) at (10,2.1) {};
\node at (10,1.55) {$R_n$};\draw[inkblue] (8.95,2.25)--(11.05,2.4);
\foreach \x in {9,9.45,9.9,10.35,10.8}{\fill[red] (\x,2.45) circle(.06);}
\node at (10.65,4.05) {读段序列};\draw[green!65!black,->,thin] (9.9,2.45)--(10.6,3.8);
\node at (1.2,4.7) {$N$};
\draw[green!65!black,->] (ch04agG)--(ch04agS);
\draw[green!65!black,->] (ch04agG)--(ch04agO);
\draw[green!65!black,->] (ch04agG)--(ch04agR);
\draw[green!65!black,->] (ch04agS)--(ch04agR);
\draw[green!65!black,->] (ch04agO)--(ch04agR);

\node[ellipse,draw=inkblue,minimum width=3cm,minimum height=1.7cm] (ch04agA) at (7.4,5.5) {};
\node[text=red,font=\ttfamily\small] at (7.4,5.85) {T C A - T G};
\node[text=green!60!black,font=\ttfamily\small] at (7.4,5.45) {T C A T T G};
\node at (7.4,5.05) {$A_n$};
\node[ellipse,draw=inkblue,minimum width=2.8cm,minimum height=1cm] (ch04agQ) at (10,.1) {$Q_n$};
\draw[green!65!black,->] (ch04agG)--(ch04agA);
\draw[green!65!black,->] (ch04agS)--(ch04agA);
\draw[green!65!black,->] (ch04agO)--(ch04agA);
\draw[green!65!black,->] (ch04agA)--(ch04agR);
\draw[green!65!black,->] (ch04agQ)--(ch04agR);

\end{tikzpicture}
\origfig{8}{允许有缺口比对的 RNA-seq 数据图形生成模型。}
\label{fig:ch04a-08}
\end{figure}


$P(G_n\mid\theta)$、$P(S_n\mid G_n)$ 和 $P(O_n\mid G_n)$ 已在第4.1.3.6.1节定义。下面计算 $P(A_n\mid G_n,S_n,O_n)$。变量 $A_n$ 表示读段 $n$ 与参考异构体 $i$ 在各比对位置的状态，状态包括匹配（M）、错配（N），以及插入／缺失或缺口（G）。从位置 $j$ 开始的读段 $R_n$ 与异构体 $G_n=i$ 之间的比对见图4.9，比对状态用 $A_n$ 表示。

设 $x$ 为比对位置，$X$ 为比对位置总数，$a[x]$ 为位置 $x$ 的比对状态，$\operatorname{trans}(a[x],a[x+1])$ 为从位置 $x$ 的状态转移到位置 $x+1$ 状态的概率。设 $\operatorname{trans}(a[X],a[X+1])$ 等于以状态 $a[X]$ 结束比对的概率，以 $\operatorname{start}(a[1])$ 表示从状态 $a[1]$ 开始的概率，则（Nariai et al. 2013）
\begin{equation}
P(A_n\mid G_n,S_n,O_n)=\operatorname{start}(a[1])\prod_{x=1}^{X}\operatorname{trans}(a[x],a[x+1]).
\tag{4.23}
\end{equation}

\begin{bookexample}{4.5}
选定的异构体 $G_n$、起始位置和读段 $n$ 见图4.9。读段 $n$ 与异构体 $G_n$ 的比对为
\[
\begin{array}{ccccccccc}
\mathrm C&\mathrm C&\mathrm A&\mathrm A&\mathrm G&\mathrm A&\mathrm G&\mathrm T&\mathrm G\\
\mathrm C&-&\mathrm C&\mathrm A&\mathrm T&\mathrm A&\mathrm G&\mathrm G&\mathrm T.
\end{array}
\]
比对状态 $A_n$ 为
\[
\mathrm{MGNMNNMNN}.
\]
图4.9给出一个简单的转移概率矩阵。假设 $\operatorname{start}(a[1])=0.5$，则
\[
\begin{aligned}
P(A_n\mid G_n,S_n,O_n)
&=0.5\times0.4\times0.5\times0.3\times0.3\times0.2\times0.3\times0.3\times0.2\\
&=3.24\times10^{-5}.
\end{aligned}
\]
\end{bookexample}
\begin{figure}[htbp]
\centering
\begin{tikzpicture}[x=1cm,y=1cm,>=Stealth,font=\small]

\node at (.4,5.4) {$G_n$};\draw[green!65!black,->] (.8,5.4)--(1.5,5.4);
\fill[red!90] (1.5,5.15) rectangle (7.2,5.65);\fill[inkblue!70] (7.2,5.15) rectangle (9.7,5.65);
\node[text=white,font=\ttfamily\small] at (4.35,5.4) {CTGAGGTTCCCAAGAGTG};
\node[text=white,font=\ttfamily\small] at (8.45,5.4) {AGGTTCA};
\draw[green!65!black] (4.6,5.7)--(4.6,6.3);\node[anchor=west] at (4.8,6.15) {起始位置 $j$};
\node at (2.2,4.5) {$R_n$};\draw[green!65!black,->] (2.6,4.5)--(3.4,4.5);
\fill[magenta!75] (3.4,4.15) rectangle (9.1,4.85);
\node[text=white,font=\ttfamily] at (6.25,4.5) {C -CATAGGT};
\draw[green!65!black,very thick,->] (5.3,4.15)--(5.3,3.55);
\node at (.9,3.1) {$A_n$};\draw[green!65!black,->] (1.3,3.1)--(1.7,3.1);
\fill[red!65!black] (1.7,2.8) rectangle (9.4,3.4);
\node[text=white,font=\ttfamily\bfseries] at (5.55,3.1) {MGNMNNMNN};
\node[anchor=west,align=left] at (9.8,3.1) {$M$：匹配\\$N$：错配\\$G$：插入／缺失};
\node[text=red,font=\bfseries,anchor=west] at (1,1.9) {转移概率矩阵};
\node[anchor=north west,inner sep=0] at (1,1.55) {\begin{tabular}{c|ccc}
&$M$&$N$&$G$\\\hline
$M$&0.5&0.3&0.4\\
$N$&0.3&0.2&0.2\\
$G$&0.2&0.5&0.4
\end{tabular}};

\end{tikzpicture}
\origfig{9}{读段与异构体之间的比对示意。}
\label{fig:ch04a-09}
\end{figure}


最后，计算给定异构体选择 $G_n=i$、方向 $O_n=k$、起始位置 $S_n=j$、比对状态 $A_n=a$ 及读段评分时，读段 $n$ 的条件概率 $P(R_n\mid G_n,O_n,S_n,A_n,Q_n)$。先为隐藏变量定义指示变量
\begin{equation}
Z_{nikja}=\begin{cases}
1,&G_n=i,\ O_n=k,\ S_n=j,\ A_n=a,\\
0,&\text{其他情形}.
\end{cases}
\tag{4.24}
\end{equation}
然后定义
\begin{equation}
\begin{aligned}
P(R_n\mid G_n,O_n,S_n,A_n,Q_n)
&=P(R_n\mid Z_{nikja},Q_n)\\
&=\prod_{x=1}^{X}\operatorname{emit}(r[x],q[x],c[x],a[x]),
\end{aligned}
\tag{4.25}
\end{equation}
其中 $x$ 为比对位置，$X$ 为比对位置总数，$\operatorname{emit}$ 为发射概率，定义为
\begin{equation}
\operatorname{emit}(r[x],q[x],c[x],a[x])=
\begin{cases}
\operatorname{subst}(r[x],q[x],c[x]),&a[x]=\text{“M”},\\
\operatorname{insert}(r[x]),&a[x]=\text{“G”},\\
\operatorname{delet}(c[x]),&a[x]=\text{“D”}.
\end{cases}
\tag{4.26}
\end{equation}
这里，“M”表示匹配／错配，“G”表示插入，“D”表示缺失。$\operatorname{subst}$ 为考虑 Phred 质量评分的替换矩阵，用于定义比对得分：在比对位置 $x$，借助质量评分 $q[x]$，为观测读段的核苷酸 $r[x]$ 与参考序列核苷酸 $c[x]$ 的比对赋分。$\operatorname{insert}$ 表示与位置无关、针对每个核苷酸定义的插入概率；$\operatorname{delet}$ 表示与位置无关、针对每个核苷酸定义的缺失概率。

同样，若读段来自噪声异构体，则由与位置无关的背景分布 $\beta$ 生成。令 $G_n=0$、起始位置 $j=1$、方向变量 $O_n=0$，来自噪声异构体的读段条件概率为
\begin{equation}
P(R_n=r_n\mid Z_{n010}=1)=\prod_{t=1}^{L}\beta(\operatorname{read}[x]),
\tag{4.27}
\end{equation}
其中 $\operatorname{read}[x]$ 表示位置 $x$ 的核苷酸，$\beta$ 的定义同前。

\bookdeepheading{paragraph}{4.1.3.6.3}{用于参数估计与 RNA-seq 数据标准化的变分贝叶斯方法}

RNA-seq 数据标准化，是给定观测读段，估计总 mRNA 中各异构体丰度所占比例，即参数向量 $\theta=[\theta_0,\theta_1,\ldots,\theta_M]^T$。换言之，最大化后验概率
\begin{equation}
P(\theta\mid R_n)=\frac{P(\theta)P(R_n\mid\theta)}{P(R_n)}.
\tag{4.28}
\end{equation}
因此，最大后验估计（MAP）为
\begin{equation}
\theta_{\mathrm{MAP}}=\argmax_{\theta}P(\theta)P(R_n\mid\theta).
\tag{4.29}
\end{equation}
然而，MAP 只考察概率密度，而不考察概率质量，可能忽略对积分或累积分布有较大贡献的部分（\url{www.cse.buffalo.edu/faculty/mbeal/thesis/beal03_2.pdf}）。

图4.6和图4.8的模型假设隐藏变量按贝叶斯网络组织。为估计参数，还需要利用隐藏变量的结构化信息。这里采用贝叶斯网络的变分方法估计参数（Nariai et al. 2013），并在此基础上引入期望最大化（expectation-maximization，EM）算法，以及变分贝叶斯 EM 算法。附录4.A说明，无缺口与有缺口比对的 EM 算法分别如下。

\par\medskip\Needspace{5\baselineskip}
\noindent\textbf{结论4.1：RNA-seq 数据无缺口比对的 EM 算法}\par
\noindent\textbf{E 步：}
\begin{equation}
\begin{aligned}
q^{(u+1)}(Z_n)
&=\frac{\dfrac{\theta_i^{(u)}}{l_i}P(R_n=r_n\mid Z_{nijk}=1,\theta^{(u)})}
{\displaystyle\sum_{i'=1}^{M}\sum_{j'=1}^{[1,\max_{i'}l_{i'}]}
\dfrac{\theta_{i'}^{(u)}}{l_{i'}}P(R_n=r_n\mid Z_{ni'j'k}=1,\theta^{(u)})},\\
&\hspace{1cm}n=1,2,\ldots,N.
\end{aligned}
\tag{4.30}
\end{equation}
以及
\par\noindent\textbf{M 步：}
\begin{equation}
\theta_i^{(u+1)}=\frac{\sum_{n=1}^{N}\sum_j\sum_kP(Z_{nijk}\mid R_n,\theta^{(u)})}{N},
\quad i=1,2,\ldots,M.
\tag{4.31}
\end{equation}

\par\medskip\Needspace{5\baselineskip}
\noindent\textbf{结论4.2：RNA-seq 数据有缺口比对的 EM 算法}\par
E 步定义为
\par\noindent\textbf{E 步：}
\begin{equation}
\begin{aligned}
q^{(u+1)}(Z_n)&=P(Z_{nijk}\mid R_n,Q_n,\theta^{(u)})\\
&=\frac{P(Z_{nijk}=1,R_n=r_n,Q_n,\theta^{(u)})}
{\sum_{(i',k',j',a')\in\pi_n}P(Z'_{ni'j'k}=1,R_n=r_n,Q_n,\theta^{(u)})},\\
&\hspace{1cm}n=1,2,\ldots,N,
\end{aligned}
\tag{4.32}
\end{equation}
其中 $\pi_n$ 是读段 $n$ 全部可能比对所对应的 $i,k,j,a$ 的集合，且 $(i,k,j,a)\in\pi_n$。
\par\Needspace{5\baselineskip}\noindent\textbf{M 步定义为：}
\begin{equation}
\theta_i^{(u+1)}=\frac{\sum_{n=1}^{N}\sum_{i,k,j,a}P(Z_{nikja}\mid R_n,Q_n,\theta^{(u)})}{N},
\quad i=1,2,\ldots,M,
\tag{4.33}
\end{equation}
其中 $(i,k,j,a)\in\pi_n$，$\pi_n$ 是读段 $n$ 全部可能比对的 $i,k,j,a$ 集合，并且
\[
\begin{aligned}
P(Z_{nikja}\mid R_n,Q_n,\theta^{(u)})
&=\frac{P(Z_{nikja}=1,R_n=r_n,Q_n,\theta^{(u)})}
{\sum_{(i',k',j',a')\in\pi_n}P(Z'_{ni'k'j'a}=1,R_n=r_n,Q_n,\theta^{(u)})},\\
&\hspace{1cm}n=1,2,\ldots,N.
\end{aligned}
\]

EM 算法的变分方法可以扩展到贝叶斯学习。用于估计无缺口及有缺口 RNA-seq 比对中转录本异构体丰度的变分贝叶斯（variational Bayesian，VB）算法，分别汇总于结论4.3与结论4.4（附录4.A；Nariai et al. 2013）。

\par\medskip\Needspace{5\baselineskip}
\noindent\textbf{结论4.3：RNA-seq 数据无缺口比对的变分贝叶斯算法}\par
\noindent\textbf{步骤1：初始化。}
对每个转录本异构体，设置 Dirichlet 分布参数的初值 $\alpha_i^{(0)}$，$i=1,\ldots,M$。

\noindent\textbf{步骤2：VBE 步。}
利用当前估计 $E_{\theta}[\theta^{(u)}]$，计算密度函数
\begin{equation}
q_Z^{(u+1)}(Z)=\prod_n\prod_i\prod_j\prod_k(\eta_{nijk})^{Z_{nijk}},
\tag{4.34}
\end{equation}
其中
\begin{equation}
\eta_{n,i,j,k}=\frac{\rho_{n,i,j,k}}
{\sum_{(i',j',k')\in\pi_n}\rho_{n',i',j',k'}}\qquad\text{或}
\tag{4.35}
\end{equation}
\[
E_Z[Z_{nijk}]=\eta_{nijk},
\]
\[
\begin{aligned}
\log\rho_{nijk}
={}&E_{\theta}[\log\theta_i]+\log P(S_n=j\mid G_n=i)\\
&+\log P(O_n=k\mid G_n=i)+\log P(R_n=r_n\mid Z_{nijk}=1),
\end{aligned}
\]
且
\[
E_{\theta}[\log\theta_i]=\psi(\alpha_i)-\psi\left(\sum_{j=1}^{M}\alpha_j\right),
\qquad
\psi(\alpha)=\frac{d\Gamma(\alpha)/d\alpha}{\Gamma(\alpha)}
\]
为双伽马函数。

\noindent\textbf{步骤3：VBM 步。}
利用当前估计 $q_Z^{(u+1)}(Z)$，计算
\begin{equation}
E_{\theta}[\theta_i^{(u+1)}]=\frac{\alpha_i^{(u+1)}}{\sum_{i=0}^{M}\alpha_i^{(u+1)}},
\tag{4.36}
\end{equation}
其中
\begin{equation}
\begin{aligned}
\alpha_i^{(u+1)}&=\alpha_i^{(u)}+\sum_{(n,j,k)\in\pi_n}E_Z[Z_{n,i,j,k}]\\
&=\alpha_i^{(u)}+\sum_{(n,j,k)\in\pi_n}\eta_{nijk}.
\end{aligned}
\tag{4.37}
\end{equation}
\noindent\textbf{步骤4：停止准则。}
若 $\lVert E_{\theta}[\theta^{(u+1)}-\theta^{(u)}]\rVert_2^2<\varepsilon$，则停止；否则返回步骤2。

\par\medskip\Needspace{5\baselineskip}
\noindent\textbf{结论4.4：RNA-seq 数据有缺口比对的变分贝叶斯算法}\par
\noindent\textbf{步骤1：初始化。}
对每个转录本异构体，设置 Dirichlet 分布参数的初值 $\alpha_i^{(0)}$，$i=1,\ldots,M$。

\noindent\textbf{步骤2：VBE 步。}
利用当前估计 $E_{\theta}[\theta^{(u)}]$，计算密度函数
\begin{equation}
q_Z^{(u+1)}(Z)=\prod_n\prod_i\prod_j\prod_k(\eta_{nijka})^{Z_{nijka}},
\tag{4.38}
\end{equation}
其中
\begin{equation}
\eta_{n,i,j,k,a}=\frac{\rho_{n,i,j,k,a}}
{\sum_{(i',j',k')\in\pi_n}\rho_{n',i',j',k',a'}}\qquad\text{或}
\tag{4.39}
\end{equation}
\[
E_Z[Z_{n,i,j,k,a}]=\eta_{n,i,j,k,a},
\]
\[
\begin{aligned}
\log\rho_{n,i,j,k,a}
={}&E_{\theta}[\log\theta_i]+\log P(S_n=j\mid G_n=i)+\log P(O_n=k\mid G_n=i),\\
&+\log P(A_n=a\mid G_n=i,S_n=j,O_n=k)\\
&+ +\log P(R_n=r_n\mid Z_{nijk}=1,Q_n),
\end{aligned}
\]
且
\[
E_{\theta}[\log\theta_i]=\psi(\alpha_i)-\psi\left(\sum_{j=1}^{M}\alpha_j\right),
\qquad
\psi(\alpha)=\frac{d\Gamma(\alpha)/d\alpha}{\Gamma(\alpha)}
\]
为双伽马函数。

\noindent\textbf{步骤3：VBM 步。}
利用当前估计 $q_Z^{(u+1)}(Z)$，计算
\begin{equation}
E_{\theta}[\theta_i^{(u+1)}]=\frac{\alpha_i^{(u+1)}}{\sum_{i=0}^{M}\alpha_i^{(u+1)}},
\tag{4.40}
\end{equation}
其中
\begin{equation}
\begin{aligned}
\alpha_i^{(u+1)}&=\alpha_i^{(u)}+\sum_{(n,j,k,a)\in\pi_n}E_Z[Z_{n,i,j,k,a}]\\
&=\alpha_i^{(u)}+\sum_{(n,j,k,a)\in\pi_n}\eta_{n,i,j,k,a}.
\end{aligned}
\tag{4.41}
\end{equation}
\noindent\textbf{步骤4：停止准则。}
若 $\lVert E_{\theta}[\theta^{(u+1)}-\theta^{(u)}]\rVert_2^2<\varepsilon$，则停止；否则返回步骤2。
\FloatBarrier
